\section{实验决策计算结果}

本节只保留用于实验选择的定量结论。默认参考点为 $P_0=300\ \mathrm{A\cdot m}$、$m_1=0.03$、$\sigma=4\ \mathrm{S/m}$、$f_s=3\ \mathrm{Hz}$。范围表使用 $P_0=\{30,100,300,1000,3000\}\ \mathrm{A\cdot m}$ 和 $m_1=\{0.01,0.03,0.1\}$。

\subsection{距离-源项-强度范围表}

表 \ref{tab:range-source-strength} 是本报告最重要的实验设计表。50--100 m 属于 near/intermediate field；300 m 以后按 far-field estimate 使用。

\begin{table}[H]
\centering
\scriptsize
\caption{距离-源项-强度范围表。静态场范围来自 $P_0$ 与观测方向扫描；轴频范围来自 $P_0$、$m_1$ 与观测方向扫描。}
\label{tab:range-source-strength}
\resizebox{\textwidth}{!}{
\begin{tabular}{cp{2.6cm}ccccc}
\toprule
距离 & 区域 & $E_{\mathrm{static}}$ 范围 (V/m) & $E_{\mathrm{static}}$ 范围 ($\mu$V/m) & $E_{\mathrm{shaft}}$ 范围 (V/m) & $E_{\mathrm{shaft}}$ 范围 ($\mu$V/m) & 基准轴频 (V/m; $\mu$V/m) \\
\midrule
30 m & near/intermediate & $2.21\times10^{-5}$--$4.42\times10^{-3}$ & $22.1$--$4421$ & $2.21\times10^{-7}$--$4.42\times10^{-4}$ & $0.221$--$442$ & $6.63\times10^{-6}$; $6.63$ \\
50 m & near/intermediate & $4.77\times10^{-6}$--$9.55\times10^{-4}$ & $4.77$--$955$ & $4.77\times10^{-8}$--$9.55\times10^{-5}$ & $0.0477$--$95.5$ & $1.43\times10^{-6}$; $1.43$ \\
100 m & near/intermediate & $5.97\times10^{-7}$--$1.19\times10^{-4}$ & $0.597$--$119$ & $5.97\times10^{-9}$--$1.19\times10^{-5}$ & $0.00597$--$11.9$ & $1.79\times10^{-7}$; $0.179$ \\
300 m & far-field estimate & $2.21\times10^{-8}$--$4.42\times10^{-6}$ & $0.0221$--$4.42$ & $2.21\times10^{-10}$--$4.42\times10^{-7}$ & $2.21\times10^{-4}$--$0.442$ & $6.63\times10^{-9}$; $0.00663$ \\
500 m & far-field estimate & $4.77\times10^{-9}$--$9.55\times10^{-7}$ & $0.00477$--$0.955$ & $4.77\times10^{-11}$--$9.55\times10^{-8}$ & $4.77\times10^{-5}$--$0.0955$ & $1.43\times10^{-9}$; $0.00143$ \\
1000 m & far-field estimate & $5.97\times10^{-10}$--$1.19\times10^{-7}$ & $5.97\times10^{-4}$--$0.119$ & $5.97\times10^{-12}$--$1.19\times10^{-8}$ & $5.97\times10^{-6}$--$0.0119$ & $1.79\times10^{-10}$; $1.79\times10^{-4}$ \\
\bottomrule
\end{tabular}}
\end{table}

直接结论：100 m 量级可以先做链路验证；300 m 量级开始考验噪声底和窄带积分；1000 m 量级的轴频基频已经进入 $10^{-10}\ \mathrm{V/m}$ 量级，必须依赖低噪声、长积分和稳定通道。

\subsection{可探测性矩阵}

可探测性使用
\[
E_{\min}(f,T_{\mathrm{int}})=\gamma S_E^{1/2}(f)/\sqrt{T_{\mathrm{int}}},
\]
其中 $\gamma=2$。判据为：信号强度高于 $E_{\min}$ 10 倍以上为“稳健可测”；高于 $E_{\min}$ 但不足 10 倍为“临界可测”；低于 $E_{\min}$ 为“不可测”。表 \ref{tab:detectability-matrix} 使用 $P_0=300\ \mathrm{A\cdot m}$、$m_1=0.03$。

\begin{table}[H]
\centering
\scriptsize
\caption{静态场与 3 Hz 轴频基频的可探测性矩阵。S100 表示 100 m 静态场，A100 表示 100 m 轴频基频，其他列同理。}
\label{tab:detectability-matrix}
\resizebox{\textwidth}{!}{
\begin{tabular}{cccccccc}
\toprule
$S_E^{1/2}$ & $T_{\mathrm{int}}$ & $E_{\min}$ & S100/A100 & S300/A300 & S1000/A1000 & 实验判断 \\
\midrule
$10^{-12}$ & 10 s & $6.32\times10^{-13}$ & 稳健/稳健 & 稳健/稳健 & 稳健/稳健 & 噪声极低，三档距离均可验证 \\
$10^{-12}$ & 100 s & $2.00\times10^{-13}$ & 稳健/稳健 & 稳健/稳健 & 稳健/稳健 & 噪声极低，三档距离均可验证 \\
$10^{-12}$ & 1000 s & $6.32\times10^{-14}$ & 稳健/稳健 & 稳健/稳健 & 稳健/稳健 & 噪声极低，三档距离均可验证 \\
$10^{-11}$ & 10 s & $6.32\times10^{-12}$ & 稳健/稳健 & 稳健/稳健 & 稳健/稳健 & 1000 m 轴频仍有余量 \\
$10^{-11}$ & 100 s & $2.00\times10^{-12}$ & 稳健/稳健 & 稳健/稳健 & 稳健/稳健 & 1000 m 轴频仍有余量 \\
$10^{-11}$ & 1000 s & $6.32\times10^{-13}$ & 稳健/稳健 & 稳健/稳健 & 稳健/稳健 & 1000 m 轴频仍有余量 \\
$10^{-10}$ & 10 s & $6.32\times10^{-11}$ & 稳健/稳健 & 稳健/稳健 & 稳健/临界 & 1000 m 轴频进入临界区 \\
$10^{-10}$ & 100 s & $2.00\times10^{-11}$ & 稳健/稳健 & 稳健/稳健 & 稳健/临界 & 1000 m 轴频进入临界区 \\
$10^{-10}$ & 1000 s & $6.32\times10^{-12}$ & 稳健/稳健 & 稳健/稳健 & 稳健/稳健 & 长积分可恢复 1000 m 轴频 \\
$10^{-9}$ & 10 s & $6.32\times10^{-10}$ & 稳健/稳健 & 稳健/稳健 & 临界/不可 & 1000 m 轴频不可作为主目标 \\
$10^{-9}$ & 100 s & $2.00\times10^{-10}$ & 稳健/稳健 & 稳健/稳健 & 稳健/不可 & 1000 m 轴频仍不足 \\
$10^{-9}$ & 1000 s & $6.32\times10^{-11}$ & 稳健/稳健 & 稳健/稳健 & 稳健/临界 & 长积分后 1000 m 轴频临界 \\
$10^{-8}$ & 10 s & $6.32\times10^{-9}$ & 稳健/稳健 & 稳健/临界 & 不可/不可 & 只适合 100 m 链路验证 \\
$10^{-8}$ & 100 s & $2.00\times10^{-9}$ & 稳健/稳健 & 稳健/临界 & 临界/不可 & 300 m 轴频临界，1000 m 不适合 \\
$10^{-8}$ & 1000 s & $6.32\times10^{-10}$ & 稳健/稳健 & 稳健/稳健 & 临界/不可 & 轴频实验仍应集中在 100--300 m \\
\bottomrule
\end{tabular}}
\end{table}

\subsection{实验等效场强设计}

桶或水槽实验不需要模拟真实舰体，只需要先模拟接收端看到的外部场强。两电极或两测点间距为 $d_{\mathrm{baseline}}$ 时，
\[
\Delta V=E_{\mathrm{target}}d_{\mathrm{baseline}}.
\]
表 \ref{tab:lab-voltage-design} 给出对应电压差，单位为 $\mu$V。换算为伏时乘以 $10^{-6}$。

\begin{table}[H]
\centering
\caption{实验等效场强设计表：$\Delta V=E_{\mathrm{target}}d_{\mathrm{baseline}}$，单位 $\mu$V。}
\label{tab:lab-voltage-design}
\begin{tabular}{cccccc}
\toprule
$E_{\mathrm{target}}$ (V/m) & 0.01 m & 0.05 m & 0.1 m & 0.5 m & 1.0 m \\
\midrule
$10^{-6}$ & 0.010 & 0.050 & 0.100 & 0.500 & 1.000 \\
$10^{-7}$ & 0.001 & 0.005 & 0.010 & 0.050 & 0.100 \\
$10^{-8}$ & $1.0\times10^{-4}$ & $5.0\times10^{-4}$ & 0.001 & 0.005 & 0.010 \\
$10^{-9}$ & $1.0\times10^{-5}$ & $5.0\times10^{-5}$ & $1.0\times10^{-4}$ & $5.0\times10^{-4}$ & 0.001 \\
\bottomrule
\end{tabular}
\end{table}

这个表直接用于判断现有 MEMS/DAQ 是否有希望测到目标量级：若系统在 $\mu$V 级以下噪声较高，$10^{-8}$--$10^{-9}\ \mathrm{V/m}$ 目标场强就需要更长基线、更低噪声或锁相/窄带积分。

\subsection{后续实验优先级}

\begin{enumerate}
    \item 第一优先：3 Hz、6 Hz、9 Hz 可控低频电偶极子注入，目标场强先做 $10^{-6}$ 到 $10^{-8}\ \mathrm{V/m}$，用于验证链路、同步、谱线检测和积分增益。
    \item 第二优先：准静态或慢变 UEP/ICCP 等效场，验证空间梯度、多通道一致性和慢漂移处理。
    \item 第三优先：50 Hz 工频干扰测量和剔除，记录每个通道的 $A_{50,i}$、相位和谐波幅值。
    \item 第四优先：尾流/流动诱导，只保留为后续单独模型，不作为当前主验证目标。
\end{enumerate}
